% ------------------------------------------------------------------------------
% 9 Discussion
% ------------------------------------------------------------------------------
\section{Discussion and Conclusions}\label{sec:discussion}

\subsection{Principal findings}

Four independent analytical paradigms, aimed at one 1980s dataset,
converge on one conclusion and enrich it from different angles. The
conclusion: levamisole plus 5-FU after resection of stage B/C colon
carcinoma delivers a large, robust treatment benefit on
disease-free survival---hazard ratio $0.62$ adjusted, $0.59$--$0.61$
Bayesian, $0.60$ under competing risks---corresponding to a
five-year absolute gain of 16.8 percentage points, a number needed
to treat of six, and 0.63 disease-free years gained in restricted
mean survival time. The enrichment: the effect is stable across seven
pre-specified subgroups with no credible interaction; a modern
group-sequential design would likely have stopped the trial for
efficacy at its first interim analysis (HR $0.52$, $Z = 3.38$ against
a boundary of $2.96$), and Bayesian monitoring would have signalled
benefit at a quarter of the evidence; covariate adjustment would have
bought real power at no inferential cost; and the largest absolute
benefits accrue to the highest-risk patients through the arithmetic of
baseline risk rather than differential response.

The re-enacted monitoring result deserves particular emphasis because
it connects historical data to contemporary trial ethics. Fifty
percent of the 506 DFS events in this dataset contain essentially all
the evidence needed for the superiority conclusion; the remaining
half of the patients' events---years of follow-up and hundreds of
recurrences---added precision but not decisions. Translated to the
1980s context in which the 625 patients randomized to observation or
levamisole alone in this very trial continued on an ineffective
strategy, an early stop would
have redirected patients to the active regimen years sooner. The
counterfactual is stylized (the historical investigators lacked
today's data-management speed and regulatory framing), but it
illustrates with unusual concreteness why alpha-spending and
beta-spending designs became the default.

\subsection{Methodological lessons}

Several results function as transferable teaching points.
\textbf{Non-collapsibility}, usually a paragraph in a textbook,
appears here as a measured quantity: with trial-calibrated prognostic
strength, the unadjusted Cox estimator targets a marginal hazard ratio
attenuated by $0.033$ on the log scale, so unadjusted and adjusted
estimates are \emph{different estimands} even in a perfectly
randomized trial (Simulation Study A). \textbf{Covariate adjustment}
raises power from 72.6\% to 80.9\% in exactly the scenario regulators
care about---modest effects, realistic sample sizes---which is the
quantitative argument behind the FDA/ICH position that adjustment
should be pre-specified and reported even when unadjusted results
match~\cite{fda2023,kahan2014}. \textbf{Competing risks} change the
estimand, not necessarily the conclusion: the naive
one-minus-Kaplan--Meier overstates cumulative recurrence by about one
percentage point at eight years, yet the Fine--Gray and cause-specific
readings agree on the treatment effect
(Section~\ref{sec:res-tx}). \textbf{Honest validation} matters more
than model sophistication: the apparent-to-cross-validated optimism
gap reaches 0.12 in the $C$-index for the most aggressive learner,
while a clinical Cox model loses nothing
(Table~\ref{tab:ml})---a reminder that in small-to-moderate clinical
datasets, disciplined regression frequently matches ensembles.
\textbf{Bayesian monitoring} demonstrates both its appeal (decisions
at 25\% of the information) and its governance question (posterior
probabilities do not self-correct for repeated looks, so their use in
confirmatory settings requires either calibration to frequentist
error rates or a genuinely decision-theoretic framing).

\subsection{Strengths}

The project's strengths are its integration and its reproducibility.
Every module runs from version-controlled code with fixed seeds,
communicates exclusively through files, and regenerates every number
and figure in this report via a single \texttt{make} invocation.
Quality-control assertions are embedded in the data pipeline (record
counts, within-patient covariate consistency, endpoint logic),
analyses are pre-specified in structure (primary composite endpoint,
adjustment set, subgroup list), and sensitivity analyses were run in
both directions---towards stronger assumptions (parametric Weibull
Bayesian model) and weaker ones (competing risks, splines, complete
case versus imputation). The simulation studies were calibrated to the
real data-generating mechanism rather than to stylized parameters, so
their conclusions transfer directly to trials of this type.

\subsection{Limitations}

Five limitations bound the interpretation. First, this is a
\emph{re-analysis of a single historical trial} from an era of
different supportive care and staging; the quantitative findings
(especially absolute risks) should not be read as contemporary
practice, and levamisole itself was later displaced by superior
regimens. Second, the subgroup and CATE analyses are exploratory:
with only 324 events in the two-arm comparison, interaction tests have limited power,
the T-learner's quartile contrasts rest on 60--100 events each, and
no external validation cohort was available to check the predicted
benefit ranking; the honest claim is methodological demonstration,
not clinical deployment. Third, the Weibull baseline is a
parametric choice; although posterior predictive checks pass and LOO
prefers it over the exponential, a piecewise-exponential or
semi-parametric Bayesian baseline would be a natural robustness
extension, as would a time-varying treatment effect in the Bayesian
model given the covariate-level PH violations detected. Fourth, the
re-enacted monitoring and Bayesian sequential analyses re-use the
final data to reconstruct interim states; the reconstruction censors
follow-up at a common time aligned to the event count and is
faithful to the ordering of the event stream, but the exercise
remains a counterfactual simulation, and early-stopping estimates
from such re-enactments are biased toward larger effects (as the
interim HR of $0.52$ versus the final $0.62$ illustrates). Fifth, the
prediction models use only ten baseline covariates from the original
case-report forms; modern depth of data (genomics, imaging) would
shift the accuracy ceiling, though the overfitting lesson of
Table~\ref{tab:ml} would only sharpen.

\subsection{Conclusions}

This capstone set out to demonstrate what a modern biostatistical
toolkit, applied with discipline to a classic randomized trial, can
extract beyond the original publication---and the answer is: the same
conclusion, reached faster (sequential design), quantified more
honestly (imputation, competing risks, absolute effects), understood
more deeply (simulation-calibrated operating characteristics,
Bayesian robustness), and personalized in direction (risk-driven
absolute benefit). The accompanying repository delivers the full
pipeline---R and Python, 22 figures, 37 tabular outputs, and this
report---reproducible with one command, in the hope that it serves
both as a portfolio piece and as a teaching resource on how
frequentist, Bayesian, and machine-learning perspectives can be made
to interrogate the same data in complementary ways.

% ------------------------------------------------------------------------------
% References
% ------------------------------------------------------------------------------
\section*{References}
\addcontentsline{toc}{section}{References}

\begin{thebibliography}{20}
\small

\bibitem{laurie1989}
Laurie JA, Moertel CG, Fleming TR, et al.
Surgical adjuvant therapy of large-bowel carcinoma: an evaluation of
levamisole and the combination of levamisole and fluorouracil.
\emph{Journal of Clinical Oncology} 1989;7(10):1447--1456. \href{https://doi.org/10.1200/JCO.1989.7.10.1447}{doi:10.1200/JCO.1989.7.10.1447}


\bibitem{moertel1990}
Moertel CG, Fleming TR, Macdonald JS, et al.
Levamisole and fluorouracil for adjuvant therapy of resected colon
carcinoma.
\emph{New England Journal of Medicine} 1990;322(6):352--358. \href{https://doi.org/10.1056/NEJM199002083220606}{doi:10.1056/NEJM199002083220606}


\bibitem{therneau2026}
Therneau TM.
\emph{A Package for Survival Analysis in R}. R package version 3.8-12,
2026. \url{https://CRAN.R-project.org/package=survival}

\bibitem{rubin1987}
Rubin DB.
\emph{Multiple Imputation for Nonresponse in Surveys}.
New York: Wiley; 1987. \href{https://doi.org/10.1002/9780470316696}{doi:10.1002/9780470316696}


\bibitem{white2009}
White IR, Royston P.
Imputing missing covariate values for the Cox model.
\emph{Statistics in Medicine} 2009;28(15):1982--1998. \href{https://doi.org/10.1002/sim.3618}{doi:10.1002/sim.3618}


\bibitem{mice2011}
van Buuren S, Groothuis-Oudshoorn K.
\emph{mice}: Multivariate Imputation by Chained Equations in R.
\emph{Journal of Statistical Software} 2011;45(3):1--67. \href{https://doi.org/10.18637/jss.v045.i03}{doi:10.18637/jss.v045.i03}


\bibitem{grambsch1994}
Grambsch PM, Therneau TM.
Proportional hazards tests and diagnostics based on weighted residuals.
\emph{Biometrika} 1994;81(3):515--526. \href{https://doi.org/10.1093/biomet/81.3.515}{doi:10.1093/biomet/81.3.515}


\bibitem{finegray1999}
Fine JP, Gray RJ.
A proportional hazards model for the subdistribution of a competing risk.
\emph{Journal of the American Statistical Association}
1999;94(446):496--509. \href{https://doi.org/10.1080/01621459.1999.10474144}{doi:10.1080/01621459.1999.10474144}


\bibitem{o1979}
O'Brien PC, Fleming TR.
A multiple testing procedure for clinical trials.
\emph{Biometrics} 1979;35(3):549--556. \href{https://doi.org/10.2307/2530245}{doi:10.2307/2530245}


\bibitem{lan-demets1983}
Lan KKG, DeMets DL.
Discrete sequential boundaries for clinical trials.
\emph{Biometrika} 1983;70(3):659--663. \href{https://doi.org/10.1093/biomet/70.3.659}{doi:10.1093/biomet/70.3.659}


\bibitem{hwang1990}
Hwang IK, Shih WJ, DeCani JS.
Group sequential designs using a family of type I error probability
spending functions.
\emph{Statistics in Medicine} 1990;9(12):1439--1445. \href{https://doi.org/10.1002/sim.4780091210}{doi:10.1002/sim.4780091210}


\bibitem{gsdesign2026}
Anderson K.
\emph{gsDesign}: Group Sequential Design. R package version 3.11.0,
2026. \url{https://CRAN.R-project.org/package=gsDesign}

\bibitem{schoenfeld1983}
Schoenfeld DA.
Sample-size formula for the proportional-hazards regression model.
\emph{Biometrics} 1983;39(2):499--503. \href{https://doi.org/10.2307/2531101}{doi:10.2307/2531101}


\bibitem{martinussen2013}
Martinussen T, Vansteelandt S.
On collapsibility and confounding bias in Cox and Aalen regression
models.
\emph{Lifetime Data Analysis} 2013;19(3):279--296. \href{https://doi.org/10.1007/s10985-013-9242-z}{doi:10.1007/s10985-013-9242-z}


\bibitem{kahan2014}
Kahan BC, Jairath V, Dor\"e CJ, Morris TP.
The risks and rewards of covariate adjustment in randomized trials:
an assessment of 12 outcomes from 8 studies.
\emph{Trials} 2014;15:139. \href{https://doi.org/10.1186/1745-6215-15-139}{doi:10.1186/1745-6215-15-139}


\bibitem{fda2023}
U.S. Food and Drug Administration.
\emph{Adjusting for Covariates in Randomized Clinical Trials for Drugs
and Biologics with Continuous Outcomes}. Guidance for Industry
(final), 2023. \href{https://www.fda.gov/regulatory-information/search-fda-guidance-documents/adjusting-covariates-randomized-clinical-trials-drugs-and-biologics-continuous-outcomes}{www.fda.gov/regulatory-information/search-fda-guidance-documents/adjusting-covariates-randomized-clinical-trials-drugs-and-biologics-continuous-outcomes}


\bibitem{andersen2003}
Andersen PK, Klein JP, Rosth{\o}j S.
Generalised linear models for correlated pseudo-observations, with
applications to multi-state models.
\emph{Biometrika} 2003;90(1):15--27. \href{https://doi.org/10.1093/biomet/90.1.15}{doi:10.1093/biomet/90.1.15}

\bibitem{parner2010}
Parner ET, Andersen PK.
Regression analysis of censored data using pseudo-observations.
\emph{The Stata Journal} 2010;10(3):408--422. \href{https://doi.org/10.1177/1536867X1001000310}{doi:10.1177/1536867X1001000310}


\bibitem{hernan2020}
Hern{\'a}n MA, Robins JM.
\emph{Causal Inference: What If}.
Boca Raton: Chapman \& Hall/CRC; 2020. \href{https://www.hsph.harvard.edu/miguel-hernan/causal-inference-book/}{www.hsph.harvard.edu/miguel-hernan/causal-inference-book/}


\bibitem{graf1999}
Graf E, Schmoor C, Sauerbrei W, Schumacher M.
Assessment and comparison of prognostic classification schemes for
survival data.
\emph{Statistics in Medicine} 1999;18(17-18):2529--2545. \href{https://pubmed.ncbi.nlm.nih.gov/10474150/}{pubmed.ncbi.nlm.nih.gov/10474150/}


\bibitem{harrell2015}
Harrell FE Jr.
\emph{Regression Modeling Strategies: With Applications to Linear
Models, Logistic and Ordinal Regression, and Survival Models}.
2nd ed. Cham: Springer; 2015. \href{https://doi.org/10.1007/978-3-319-19425-7}{doi:10.1007/978-3-319-19425-7}


\bibitem{gelman2013}
Gelman A, Carlin JB, Stern HS, Dunson DB, Vehtari A, Rubin DB.
\emph{Bayesian Data Analysis}. 3rd ed. Boca Raton: CRC Press; 2013. \href{https://doi.org/10.1201/b16618}{doi:10.1201/b16618}


\bibitem{salvatier2016}
Salvatier J, Wiecki TV, Fonnesbeck C.
Probabilistic programming in Python using PyMC3.
\emph{PeerJ Computer Science} 2016;2:e55.
(PyMC 6, \url{https://www.pymc.io}). \href{https://doi.org/10.7717/peerj-cs.55}{doi:10.7717/peerj-cs.55}


\bibitem{vehtari2017}
Vehtari A, Gelman A, Gabry J.
Practical Bayesian model evaluation using leave-one-out cross-validation
and WAIC.
\emph{Statistics and Computing} 2017;27(5):1413--1432. \href{https://doi.org/10.1007/s11222-016-9696-4}{doi:10.1007/s11222-016-9696-4}


\bibitem{royston2011}
Royston P, Parmar MKB.
The use of restricted mean survival time to estimate the treatment
effect in randomized clinical trials when the proportional hazards
assumption is in doubt.
\emph{Statistics in Medicine} 2011;30(20):2409--2421. \href{https://doi.org/10.1002/sim.4274}{doi:10.1002/sim.4274}


\bibitem{uno2014}
Uno H, Claggett B, Tian L, et al.
Moving beyond the hazard ratio in quantifying the between-group
difference in survival analysis.
\emph{Journal of Clinical Oncology} 2014;32(22):2380--2385. \href{https://doi.org/10.1200/jco.2014.55.2208}{doi:10.1200/jco.2014.55.2208}

\end{thebibliography}

